%% bare_conf.tex    
%% V1.4b
%% 2015/08/26
%% by Michael Shell
%% See:
%% http://www.michaelshell.org/
%% for current contact information.
%%
%% This is a skeleton file demonstrating the use of IEEEtran.cls
%% (requires IEEEtran.cls version 1.8b or later) with an IEEE
%% conference paper.
%%
%% Support sites:
%% http://www.michaelshell.org/tex/ieeetran/
%% http://www.ctan.org/pkg/ieeetran
%% and
%% http://www.ieee.org/

%\documentclass[usletter, 10 pt, conference]{IEEEtran}
\documentclass[a4paper, 10 pt, conference]{IEEEtran}
\usepackage{eurosym}
\usepackage{upgreek}
\usepackage[left=1.9cm, right=1.3cm, bottom=2.1cm, top=3.66cm]{geometry}

\usepackage{cite}

% *** GRAPHICS RELATED PACKAGES ***
%
\ifCLASSINFOpdf
  \usepackage[pdftex]{graphicx}
  \usepackage{import}

% correct bad hyphenation here
\hyphenation{op-tical net-works semi-conduc-tor}
%\usepackage{anyfontsize}
\usepackage[english]{babel} 
\usepackage[utf8]{inputenc}
\usepackage[]{caption}
\captionsetup{font=small}   
\usepackage{subcaption}
\usepackage{placeins}
\usepackage{color}
\usepackage{mathtools}
\usepackage{placeins}
%\usepackage{graphicx,import}
\usepackage{relsize}
\usepackage{amssymb}
\usepackage{units}
\usepackage{amsmath}
\usepackage{float} 		
\renewcommand{\figurename}{Fig.}
\addto\captionsenglish{\renewcommand{\figurename}{Fig.}}
\usepackage{nicefrac}
\usepackage{tikz}
\usepackage[percent]{overpic}

%\usepackage[table]{xcolor}
\setlength{\arrayrulewidth}{1mm}
%\arrayrulecolor[HTML]{DB5800}
%\usepackage{ragged2e}
\definecolor{mygray}{RGB}{70, 70, 70}
\definecolor{mygray2}{RGB}{40, 40, 40}

\usepackage{pgf}
\usepackage{tikz}
\usetikzlibrary{arrows,automata}
\usetikzlibrary{positioning}
\tikzset{    state2/.style={ rectangle,  draw=black, inner sep=7pt,  text centered  },}
\usepackage[color=black,opacity=1,angle=0,scale=1]{background}
\backgroundsetup{
% contents={\textcopyright 2016 IEEE. Personal use of this material is permitted. Permission from IEEE must be obtained for all other uses in any current or future media \newline, including reprinting/republishing this material for advertising or promotional purposes, creating new collective works, for resale or redistribution to servers or lists, or reuse of any copyrighted component of this work in other works.},
contents={
\begin{tikzpicture}
\node at (current page.center) [align=center] {\textcopyright 2020 IEEE. Personal use of this material is permitted. 
\\ Permission from IEEE must be obtained for all other uses, in any current or future media, including reprinting/republishing this material for advertising or promotional purposes, 
\\ creating new collective works, for resale or redistribution to servers or lists, or reuse of any copyrighted component of this work in other works.
\\ DOI: 10.1109/ITSC45102.2020.9294452};
\end{tikzpicture}},
placement=bottom,
scale=0.6,
vshift=20
}


\begin{document}
\title{\vspace*{0.4cm} 
Proactive Risk Navigation System for Real-World Urban Intersections
\vspace*{-0.3cm}} %complex
\author{\IEEEauthorblockN{Tim Puphal$^{1\star}$, Benedict Flade$^{1\star}$, Daan de Geus$^{2}$ and Julian Eggert$^{1}$} %ast
\IEEEauthorblockA{$^{1}$ Honda Research Institute (HRI) Europe, 
	Carl-Legien-Str. 30, 63073 Offenbach, Germany \\
	Email: {\tt\small \{tim.puphal, benedict.flade, julian.eggert\}@honda-ri.de} \\
	$^{2}$ Eindhoven University of Technology (TU/e), 5612 AZ Eindhoven, the Netherlands \\
	Email: {\tt\small d.c.d.geus@tue.nl}} %\{d.c.d.geus, g.dubbelman\}@tue.nl}} 
    $\star$ The authors contributed equally to this work}
    
% make the title area
\maketitle

% maximal length 8 pages
% As a general rule, do not put math, special symbols or citations in the abstract
\begin{abstract} 
We consider the problem of intelligently navigating through complex traffic. Urban situations are defined by the underlying map structure and special regulatory objects of e.g. a stop line or crosswalk. Thereon dynamic vehicles (cars, bicycles, etc.) move forward, while trying to keep accident risks low. 

Especially at intersections, the combination and interaction of traffic elements is diverse and human drivers need to focus on specific elements which are critical for their behavior.
To support the analysis, we present in this paper the so-called Risk Navigation System (RNS). RNS leverages a graph-based local dynamic map with Time-To-X indicators for extracting upcoming sharp curves, intersection zones and possible vehicle-to-object collision points. 

In real car recordings, recommended velocity profiles to avoid risks are visualized within a 2D environment. 
By focusing on communicating not only the positional but also the temporal relation, RNS potentially helps to enhance awareness and prediction capabilities of the user. 
\end{abstract}

% For peerreview papers, this IEEEtran command inserts a page break and
% creates the second title. It will be ignored for other modes.
\IEEEpeerreviewmaketitle

\section{Introduction}

The increasing complexity of source code poses a key challenge to the reliability of large-scale software systems. Software bugs in these systems can lead to safety issues~\cite{bug_safety} for users around the world as well as cause non-negligible financial losses~\cite{bug_loss}. As such, developers have to spend a large amount of time and effort on bug fixing. Consequently, \aprfull (\apr), designed to automatically generate patches to fix software bugs, has attracted wide attention from both academia and industry~\cite{long2016prophet, legoues2012genprog, long2015spr, lou2020can, tufano2018empstudy}. 


To achieve \apr, one popular approach is known as Generate-and-Validate (G\&V)~\cite{qi2015gv, ghanbari2019prapr, lou2020can, le2016hdrepair, legoues2012genprog, wen2018capgen, hua2018sketchfix, martinez2016astor, koyuncu2020fixminder, liu2019tbar, liu2019avatar}, which is typically based on the following pipeline: First, fault localization techniques~\cite{wong2016fl, abreu2007ochiai, zhang2013injecting, papadakis2015metallaxis, li2019deepfl, li2017transforming} are applied to determine the suspicious locations in programs where bugs are likely to exist. Then, the buggy locations are used by the \apr tools to generate a list of patches that replace buggy lines with correct lines. Afterward, each patch is validated against the original test suite to identify any \emph{plausible patches} (i.e., passing all tests in the test suite). Finally, to determine the \emph{correct patches}, developers examine the list of plausible patches to see if any of them can correctly fix the bug. 

Traditional \apr tools can mainly be categorized into heuristic-based~\cite{legoues2012genprog, le2016hdrepair, wen2018capgen}, constraint-based~\cite{mechtaev2016angelix, le2017s3, demacro2014nopol, long2015spr} and \template~\cite{ghanbari2019prapr, hua2018sketchfix, martinez2016astor, liu2019tbar, liu2019avatar}. Among these traditional tools, \template \apr tools~\cite{ghanbari2019prapr, liu2019tbar, benton2020effectiveness} have been able to achieve state-of-the-art results. \Template \apr tools typically leverage pre-defined templates (e.g., adding a nullness check) for bug fixing. However, since these fix templates are typically handcrafted, the number and types of bugs they are able to fix can be limited. 



To address the limitations of traditional \apr, researchers have proposed various \learning \apr tools~\cite{li2020dlfix, chen2018sequencer, jiang2021cure, lutellier2020coconut, zhu2021recoder, ye2022rewardrepair} based on the \nmtfull (\nmt) architecture~\cite{sutskever2014mt} where the input is the buggy code snippets and the goal is to translate the buggy code snippets into a fixed version. To accomplish this, \learning \apr tools require supervised training datasets with pairs of both buggy and fixed code snippets in order to learn how to perform this translation step. These training data are usually obtained by mining historical bug fixes using heuristics/keywords~\cite{dallmeier2007benchmark}, which can be imprecise for identifying bug-fixing commits; even the actual bug-fixing commits can include irrelevant code changes, leading to further pollution in the dataset~\cite{xia2022alpharepair}.
% 
Moreover, it can be hard for such \apr tools to generalize and fix bug types unseen during training. 



To better leverage recent advances in \plmfull{s} (\plm{s}), researchers~\cite{xia2022alpharepair, xia2023repairstudy, kolak2022patch, prenner2021codexws} have directly applied \plm{s} to generate patches without bug-fixing datasets. These \llm-based \apr tools work by either directly generating a complete code function~\cite{prenner2021codexws, xia2023repairstudy} or predict/infill the correct code snippet given its surrounding context~\cite{xia2022alpharepair, xia2023repairstudy}. By directly using \llm{s} that are pre-trained on billions of open-source code snippets, \llm-based \apr tools can achieve state-of-the-art performance on many repair datasets~\cite{xia2022alpharepair}. 


% 
%
%

Traditional \apr tools have long used the insight of the \emph{plastic surgery hypothesis}~\cite{barr2014plastic} where it states that the code ingredients to fix a bug already exist within the same project. Traditional \apr tools have manually designed pattern-~\cite{ghanbari2019prapr, saha2017elixir} or heuristic-based~\cite{jiang2018simfix, legoues2012genprog} approaches to finding and using such relevant code ingredients to generate fixes for bugs. However, the plastic surgery hypothesis has been largely ignored in \llm-based \apr. In fact, \llm provides a unique opportunity to fully automate the plastic surgery hypothesis idea via fine-tuning (learning project-specific information via model updates from the buggy project) and prompting (directly providing relevant code ingredients to the model), and make it directly applicable to different languages (since the \llm{s} are typically multi-lingual).%
Moreover, despite the intensive manual efforts involved, traditional \apr tools still cannot fully leverage project-specific information due to large search space for leveraging/composing existing code ingredients. In contrast, the project-specific information can effectively leveraged by \llm{s} due to their power in code understanding/vectorization, e.g., even partial/imprecise information may still guide \llm{s} in correct patch generation!
 To this end, we ask the question: \emph{How useful is the plastic surgery hypothesis in the era of \plm{s}}?








\mypara{Our Work.} To answer the question, we present \ourtech{\xspace} -- a \llm-based approach that automatically utilizes the plastic surgery hypothesis by systematically combining multiple fine-tuning and prompting strategies for \apr. \ourtech fine-tunes \plm{s} using two novel domain-specific training strategies: \textbf{\epfinetune} -- we fine-tune using the original buggy project by aggressively masking out a high percentage of tokens, which allows \plm to learn project-specific code tokens and programming styles; and \textbf{\rofinetune} -- which only masks out a single continuous code sequence per training sample, allowing the model to get used to the final \csapr task of predicting a single continuous code sequence. Furthermore, we directly leverage the ability for \plm{s} to understand natural language instructions and introduce a novel prompting strategy, \textbf{\idprompting}, which uses information retrieval and static analysis to obtain a list of relevant identifiers for the buggy lines. While such relevant identifiers are critical for fixing some difficult bugs, they may not be seen by the \llm during inference due to limited context window size. Through the use of prompting, we directly tell the model to use these extracted identifiers (relevant code ingredients) to generate the correct code. Finally, to perform repair, we combine all four model variants (including the base model, both fine-tuned models and the base model with prompting) for the final repair.





While our insight of leveraging the plastic surgery hypothesis for \llm-based \apr is generalizable across different types of \plm{s}, to implement \ourtech, we choose a recent \plm{\xspace}, \ctfive~\cite{wang2021codet5}, which is pre-trained on millions of open-source code snippets. \ctfive is an encoder-decoder model trained using \mspfull (\msp) objective where a percentage of tokens are masked out and each continuous masked token sequence is referred to as a masked span. Also, although we only extract relevant identifiers from the current buggy project (since this paper focuses on the plastic surgery hypothesis), our work can be easily extended to obtain other code information (such as relevant statements or functions) from other sources, such as  the massive pre-training corpora~\cite{husain2020codesearchnet} or historical bug-fixing datasets~\cite{jiang2019infer}, which can provide more coding knowledge for \llm{s}. Besides, although we mainly focus on using traditional string comparison algorithms for information retrieval in this paper, these techniques can be easily replaced by other frequency-based retrieval~\cite{robertson2009probabilistic} and neural search (or embedding-based search)~\cite{reimers2019sentence}.
  In summary, this paper makes the following contributions:


%


\begin{itemize}[noitemsep, leftmargin=*, topsep=0pt]
    \item \textbf{Dimension.} This paper is the first to revisit the important plastic surgery hypothesis in the era of \llm{s}. It opens up a new dimension for \llm-based \apr to incorporate previously neglected information from the buggy project itself to boost \apr performance. Furthermore, it demonstrates the promising future of retrieval-based prompting for modern \llm-based \apr.
    \item \textbf{Implementation.} We implement \ourtech based on the recent \ctfive model. We augment the model using two novel fine-tuning strategies: \epfinetune and \rofinetune, along with a novel prompting strategy based on information retrieval and static analysis: \idprompting. We combine the patches generated by all four models together and perform patch ranking to speed up \apr.% 
    \item \textbf{Evaluation Study.} We conduct an extensive evaluation against state-of-the-art \apr tools. On the widely studied \dfj 1.2 and 2.0 datasets~\cite{just2014dfj}, \ourtech is able to achieve the new state-of-the-art results of 89 and 44 correct bug fixes (15 and 8 more than best baseline) respectively.  Furthermore, we perform a broad ablation study to justify our design. \ourtech demonstrates for the first time that the plastic surgery hypothesis can substantially boost \llm-based \apr and advance state-of-the-art \apr, while being fully automated and general. Moreover, even partial/imprecise code ingredients may still effectively guide \llm{s} for \apr!
\end{itemize}



\section{Map-Based Risk Prediction}
\label{sec:riskextr}
%Representation of urban situations requires a suitable environment model. 
The target is to represent driving situations with a flexible and efficient environment model. Generalized ADAS for intersections require clear data management functions.  %Representation of driving situations requires a flexible environment model. 
The following section \ref{subsec:ldm} shows how a graph-based local dynamic map can be leveraged for this purpose. Additionally, risk-relevant data within a certain virtual horizon is retrieved to allow predictive trajectory evaluations and planning in section \ref{subsec:trajeval}. %We combine simple but diverse 

\subsection{Relational Local Dynamic Map (R-LDM)}
\label{subsec:ldm}
Underlying road geometry significantly influences and constrains all traffic participants. %in urban cities multi-faceted  is available for the ego driver.  %Additionally, there are many different informations available. 
From the ego driver, multi-faceted sensory inputs have to be processed in real-time.
Since important meta data lies in the connectivity of environment entities,
a graph representation, consiting of nodes, attributes and relations, is beneficial. 
% In this sense, it has items as nodes and more importantly attributes as connections. 
% This also yields a natural link to road structure data, since the road map interconnections and geometry data build a graph. 
The Relational Local Dynamic Map (R-LDM) presents an abstraction layer between sensor data and higher level functions~\cite{ldm2017}.
Such functions include e.g. risk prediction and visualization as required by our RNS. 
With the R-LDM, we support the paradigm shift from single sensor-controller loops towards technologies with enhanced understanding.

\subsubsection{Graph Structure}

A R-LDM graph stores data on four layers based on their dynamicity (compare Fig. \ref{fig:rldm} top right).
The map is represented as nodes with additional data of shape, orientation or type saved as attributes. % at intersections, lane marking types and other information like height and curvature.
On the lowest layer, the street geometry presents one of the most crucial data types. Generally, \textbf{static data} is stored in three levels of granularity, ranging from detailed lanes via half roads to the whole road.
In this context, half roads are the sum of all lanes pointing in the same direction.
Furthermore, the road subgraph consists of alternating segments and junctions, while all junctions at one location forms an intersection.

The second layer comprises \textbf{quasi-static data} with changes on the scale of several days. 
This implies traffic lights, roadside infrastructure, traffic signs or construction sites. 
% The mentioned data falls into static map data which is expected to change on a very slow timescale and quasi-static data with changes on a scale of several days and longer. Here, we have besides traffic lights, also roadside infrastructure, traffic signs or construction sites. 
On the third layer, \textbf{transient data} is stored changing on hours scale or faster (i.e., traffic light states, car density or slippery road conditions). 
Finally, the fourth layer contains highly \textbf{dynamic data} about the state of traffic participants, including vehicles and pedestrians. 
Detected dynamic objects are connected to specific nodes of the lower layers, which can be used for risk extraction in the ego vicinity.
The R-LDM is consequently maintained or updated in real-time.

%, which allow further behavior prediction, planning and support.
With regard to lower layers, available map data sources such as the crowd-sourced OpenStreetMap (OSM)\footnote{See www.openstreetmap.org for further explanation.} can be applied. % \cite{OSM}. %~\cite{OSM}.
% While the lowest layer relies mostly on offline data from OSM, it is the highest layer which comes from post-processed sensor information. 
%STATIC DATA INPUT
% For the download of the map data, we rely on OpenStreetMap (OSM), a crowd-sourced database that is freely available to the public \cite{OSM}. 
OSM mainly consists of navigational information represented in the form of road centerlines and intersection points. 
Therefore, we add topological information on lane level and estimate the actual lane geometry under assumption of a concrete lane width. 
% Afterwards, by combining the lane semantics with an estimate of the lane and road width, we obtain the actual geometric road shape. 
Besides, stop lines are manually added and linked to incoming lanes, while the positions of traffic lights, crosswalks or buildings are similarly extracted from OSM. 

\vspace{0.06cm}
\subsubsection{Virtual Horizon}
\label{sub:horizon}

The graph structure with its connected nodes suits itself to a wide scope of potential queries with regard to path finding, routing or tree search. 
In this way, extraction of behaviorally relevant situations is possible in a straightforward way.
One crucial information for risk estimation is the knowledge of the upcoming driving path of traffic participants.
When considering start and end points in the map, we can intuitively query a \textit{sequence of nodes} (segments and junctions) along the route by traversing its relations. 

Fig. \ref{fig:rldm} shows an example of an intersection scenario. 
Therein, we obtain lane nodes (:LaneSegment, :LaneJunction) and traffic light nodes (:TrafficLight), while the graph is explored based on the current position and orientation of a chosen vehicle (:Vehicle). 
Relevant data is stored in the attributes of the individual nodes. 
Essentially, we can query any processed information that has been acquired online (e.g. from sensor data) or offline (e.g. from OSM).

The left part of Fig. \ref{fig:rldm} depicts the corresponding excerpt of the graph concept that highlights layers (:StaticEntities, :QuasiStaticEntities, :TransientEntities, :DynamicEntities) and granularity.
In order to successfully calculate risk, not only one but several \textit{possible paths} have to be queried. 
In particular, a virtual horizon can be determined which extracts predicted paths for the vehicles in the traffic scene.

\begin{figure}
  \centering
  \resizebox{\linewidth}{!}{
  \import{img/}{rldm.pdf_tex}}
  \caption{Nodes hierarchy and interconnection graph structure within R-LDM for an intersection with traffic lights.}
  \label{fig:rldm}
\end{figure}
\vspace{-0.05cm}

\vspace{0.05cm}
To illustrate this virtual horizon, Fig. \ref{fig:virtualhorizon} addresses a scenario with two approaching cars. 
First, we map the GNSS-based ego vehicle position (in green) onto the center of the corresponding lane element.
% This can be achieved, because the OSM information is given in world coordinates. In this process, we also use . 
Afterwards, we retrieve other cars from lidar (in red) up to a certain distance with R-LDM queries. More stored elements, for instance, quasi-static data with crosswalks can be equally obtained.

By scanning lane nodes for all cars, concatenating their polylines and ordering the branches (straight, left and right, 
etc.) according to a \textit{tree structure}, we obtain distinct paths. In Fig. \ref{fig:virtualhorizon}, every car posseses three possible paths.
The root of the tree is the actual position and the length increases with the number of segments we traverse.
With the virtual horizon, we can directly estimate simple risks. Crossings of the paths reflect conflict zones, where accidents might happen. These are unique and critical for the given road.

\subsection{Trajectory Evaluation}
\label{subsec:trajeval}
After querying the R-LDM for upcoming map structures, traffic elements and sensed objects, we predict trajectories from the vehicle dynamics and evaluate their concrete risk time-courses. Proper safe maneuvers or velocities for recommendation are related to our behavior planning. We divide risks for this purpose into three generic types: 1) collision risk, 2) curvature risk and 3) regulatory risk. In RNS, we combine the risk types to holistically support the driver.
%-Risk sources vs actual risk \\
%-Collision risk (next intersection, intersection zone, intersecting possible paths or intersection lanes) $\rightarrow$ from LDM queries \\

\begin{figure*}
  \centering
  \resizebox{0.82\linewidth}{!}{
  \import{img/}{vhorizon.pdf_tex}}
  \caption{Steps of the virtual horizon shown exemplarily for collision risk. Left: Filtered relevant cars and their positions. Middle: Extraction of possible future paths. Right: Map-based risk extraction, with crosses indicating future path intersections.}
  \label{fig:virtualhorizon} 
\end{figure*} 

\begin{figure}
  \centering
  \resizebox{0.80\linewidth}{!}{
  \import{img/}{type_of_risks.pdf_tex}}
  \vspace{0.08cm}
  \caption{Three risk types of RNS with time as well as space metrics (collision, curve and regulatory risks).}
  \label{fig:risktypes} 
\end{figure} 

\vspace{0.06cm}
\subsubsection{Collision Risk} %\ref{sub:horizon}
Besides coarse conflict zones from the last section \ref{subsec:ldm}, risk metrics of quantitative collision probability can be obtained as well.
In an initial step, paths are transformed with constant velocity models into trajectories. For a car-to-car encounter (indexed 1,2), this gives us position sequences $\mathbf{x}_1(s)$ and $\mathbf{x}_2(s)$ over the predicted time $s$ with assumed constant velocities $v_1$ and $v_2$. 

The distance between the vehicles represents the pointwise trajectory difference and can be written as 
\begin{equation}
d(s)=\left\| \mathbf{x}_2(s) - \mathbf{x}_1(s) \right\|.
\end{equation}
For time risks, we next consider the time of maximal criticality along the predicted trajectories. In other words, we filter out the minimal distance, which is 
the event distance $d_E$ and get the corresponding Time-To-Closest-Encounter (TTCE) as $s_E$. The two indicators are in summary retrieved with

\begin{equation}
d_E=\text{min}\{d(s)\} \text{ and } s_E=\text{argmin}\{d(s)\}.
\end{equation}
\vspace{0.05cm}
\hspace{-0.114cm}Coarsely, collision risk corresponds to the inverse $1/s_E$ and is only considered if $d_E \hspace{-0.02cm} < \hspace{-0.02cm} d_{\text{min}}$ holds true. The minimal value $d_{\text{min}}$ defines the risk sensitivity. 

Once there are multiple trajectory pairs as in Fig. \ref{fig:virtualhorizon}, we would iterate over the combinations $\text{\textbf{x}}_{ij}$ (i.e., variation of ego paths $i$ and other paths $j$) and extract the highest risk by $\text{max}\{1/s_{E,ij}\}$. In this method, we therefore capture time and space of passing vehicles. %\textcolor{red}{emphasize is basic risk, but can also use more 

\subsubsection{Curvature Risk}
From quantitative risks, we can also infer behavior recommendations as target velocities that minimize criticality and maximize ego utility. This is essentially a harder problem, since the driver should now understand the warning as well as interpret and process the advice. 

Consider the situation that the ego user is driving with velocity $v_0$ at the current time $t$. We define a turn segment with curvatures $[\kappa_{\text{start}}, \kappa_{\text{end}}]$ in the extracted path, if a thresholding condition
\begin{equation}
\kappa_{\text{start}} > \kappa_{\text{th}} \text{ and } \kappa_{\text{end}} < \kappa_{\text{th}}
\end{equation}
is valid. Here, the threshold $\kappa_{\text{th}}$, similarly as $d_{\text{min}}$, lets us adjust the turn detection along the future profile $\kappa(s)$. 

In curves, the vehicle is prone to drive off the lane. For higher velocities $v_0$, increased lateral acceleration is exerted. When $a_y$ depicts the dynamics limit that the vehicle can follow and $\kappa_{\text{max}} \hspace{-0.07cm} = \hspace{-0.07cm} \text{max}\{\kappa(s)\}$ is the maximal curvature in the segment, we can derive a target velocity via
\begin{equation}
v_{\text{tar}} = \sqrt{a_y/\kappa_{\text{max}}}.
\end{equation}
The variable $v_{\text{tar}}$ indicates the speed that should be maximally reached at the curve.  For low risk, we should accordingly move slower than $v_{\text{tar}}$, which leads to the condition $v_0 < v_{\text{tar}}$. Driving with $v_0 \ll v_{\text{tar}}$ is also not recommended, because of the reduction in utility.

%Curve risk is extracted solely based on the road. Another structural risk that we could retrieve with the R-LDM is occlusion risk. The visibility of drivers decreases around buildings and relevant cars are overlooked.

\subsubsection{Regulatory Risk}
Quasi-static elements of a stop line, traffic light or crosswalk create risks at rule violation. The reason behind this circumstance is that generally not obeying the norms can lead to unexpected situations with accidents. Only because the vehicles are rule-conform, we can drive safely. %and estimate the intentions of other vehicles.

In this sense, we may want to stop directly in front of the traffic data. Mathematically, we assume a soft braking trajectory with constant deceleration $a_c$ to the stopping point in the distance $d_c$ along the ego path. With the time $s$, we initially obtain the basic kinematic equation
%stopping distance ds and dc
\begin{equation}
d = \frac{a \cdot s^2}{2}.
\label{eq:distancestop}
\end{equation}
Substituting $s$ with $v_{\text{tar}}/a_c$ and inserting $a=a_c$ as well as $d=d_c$ in Eq. (\ref{eq:distancestop}) ultimately results into
\begin{equation}
v_{\text{tar}} = \sqrt{2a_c \cdot d_c},
\label{eq:vtar}
\end{equation}
whereby $v_{\text{tar}}$ is the target velocity to be able to stop at the intersection with certainty.\footnote{For Eq. (\ref{eq:vtar}), we do not account for the reaction time $t_r$. However, $t_r$ can be added for large distances $d_c$ with $v_{\text{tar}} \approx \sqrt{2a_c \cdot d_c} - a_c\cdot t_r$.} 
%wolframalpha https://www.wolframalpha.com/input/?i=%5Bsqrt%282x%29%2C+sqrt%282x%29*sqrt%281%2B%280.8%5E2%29%2F%282*x%29%29-0.8%2C+sqrt%282x%29-0.8%5D+from+0+to+10

Fig. \ref{fig:risktypes} illustrates the corresponding variables for the three types of risk from this section. We want to stress at this point that our approach employs explicit trajectory prediction along the map paths from the R-LDM. Concretely, we assume constant velocity for the trajectory prediction in the collision case, an acceleration or deceleration trajectory to a fixed reduced velocity in the curve case and a smooth braking trajectory to a fixed position for the regulatory case. Therewith, we eventually filter and consider single time points in the trajectory for the risks.

\section{Visualization On Demand} %Visualization Elements
\label{sec:visrisk}
Based on environment data and trajectory evaluation, we now present ways of communicating the situation and risks on a visual display to achieve an ADAS.
In this context, we employ a renderer that visualizes all the information in a joint Cartesian coordinate system (see section \ref{subsec:sim}). 
Once driving risks are detected, design elements are overlayed on the display with section \ref{subsec:active} and section \ref{subsec:warning}. 

\subsection{Simulator Environment}
\label{subsec:sim}
Nodes of the R-LDM have a range of potential attributes, such as the 3D position or geometrical shape of objects. 
% For instance, the road centerline is a polyline with bounderies to the left and right. Crosswalks have a defined width and buildings a polygonal outline description. 
In the renderer, we always visualize static and quasi-static data that lie in the field of view from the ego vehicle. 
For this, a local 3D model is generated by converting geographic points with (lat, lon, alt) into Cartesian coordinates of (x, y, z). 
% and project the positonal relations from a view perspective with a transformation matrix. 
Fig. \ref{fig:3Dsimulator} depicts an exemplary map section having several intersections in bird's-eye view.
% with several intersections, stop lines and crosswalks. 
On the top right, the first person view of a vehicle approaching a crosswalk is shown. 

The dynamic data is then added to this static view. A zoomed-in excerpt from the map is given at the bottom of Fig. \ref{fig:3Dsimulator} that includes a recorded GNSS trace (red).
We project the trace onto the connected lane center, which is pictured in green. 
% Because we project the ego position on the closest lane segment, on the bottom right the measured trace is changed in red and the aligned trace is marked in green.
Consequently, the virtual horizon and its possible paths are retrieved as described in section \ref{subsec:ldm}. 
We can lastly update and move the excerpt with the current position from the GNSS to obtain a live simulation.

\subsection{Proactive Support}
\label{subsec:active}
Communication of spatial as well as spatio-temporal relations is crucial for risk-averse driver support. 
% This has the reason that humans can estimate the time better than positions (especially for risks). 
% Velocity contains implicitly the time as well. 
Further sources of information are cause, likelihood and severity of a potential risks.  
% if a collision happens. 
The next step for RNS is the choice of suitable design elements. 
In this process, we suppose that we know where the ego vehicle is driving (i.e., the ego path) from its navigation route. 
Yet, for surrounding vehicles, all paths are considered.

\subsubsection{Hazard Route Element}
The so-called hazard route in Fig. \ref{fig:charts} is a concept that consists of a scale portraying distances to an upcoming risk element.
Furthermore, the geometrical area or length of risks is considered.
Risk is thus measured with respect to the ego path, ranging from the current position  $\Delta l \hspace{-0.03cm}=\hspace{-0.03cm} \unit[0]{m}$ to the end of the path $\Delta l_{h}$.
Here, the length $\Delta l_{h}$ can be chosen according to own preferences. 

At an upcoming intersection, risk is defined by the section of the path that lies within the junction.
Since risk corresponds to exposition time, we encode the path part from the intersection $I_z$ with a color, ranging from green for short intersections to red for long ones. 
%allgemein risiko entlang des pfades zu intersection zone
%share of junction segment to navigation route + 
%one case with large intersection far and one case with small intersection close
Fig. \ref{fig:charts}~a) gives two examples of the hazard route.
The left bar shows a large intersection (e.g. multi-lane four-way stop) in vicinity and the right bar has a small and consecutive medium junction. 
% In the case of collision risk, the intersection zone $I_z$ can be used.
% Depending on the value of $I_z$ (low, medium and large), the area is marked from green, to yellow until red for conveying the criticality. 
This emphasizes that we may include more than one intersection in our warnings.

\begin{figure}[t]
  \centering
  \includegraphics[width=0.95\linewidth]{./img/simulator.png}
  \caption{Rendered road network from two perspectives with the ego position being projected on the navigation route. \vspace{0.45cm}}
  \label{fig:3Dsimulator}
\end{figure}

\begin{figure}[t]
  \centering
  \resizebox{\linewidth}{!}{
  \import{img/}{velocity_scale_new.pdf_tex}}  
  \caption{Chart elements for proactive support. Hazard route (left) and velocity scale (right).} %\vspace{-0.3cm}}
  \label{fig:charts} 
\end{figure} 

\subsubsection{Velocity Scale Element}
The velocity scale, Fig. \ref{fig:charts}~b), is a second chart element which qualifies the difference between the current velocity of the vehicle $v_0$ and the target velocity $v_{\text{tar}}$ from the trajectory evaluation of section \ref{subsec:trajeval}. 
The scale shows possible velocity values, from standstill $v\hspace{-0.05cm}=\hspace{-0.05cm}\unit[0]{m/s}$ to a maximal velocity $v_{\text{max}}$. Depending on the difference $|v_0 \hspace{0.05cm} - \hspace{0.05cm} v_{\text{tar}}|$, the situation is rated as safe with $v_0 \hspace{-0.042cm} \approx \hspace{-0.042cm} v_{\text{tar}}$ (green, left), as dangerous with e.g. $v_0 \hspace{-0.05cm} < \hspace{-0.05cm} v_{\text{tar}}$ (yellow, middle) to critical with $v_0 \hspace{-0.07cm} \ll \hspace{-0.07cm} v_{\text{tar}}$ (red, right). The same cases hold true for the opposite circumstances, i.e., $v_0 \hspace{-0.032cm} > \hspace{-0.032cm} v_{\text{tar}}$. 
This velocity scale can be employed for curve or regulatory risks. 
Moreover, we may set an enforced speed limit as the target velocity $v_{\text{tar}}$ for proactive behavior, once there is no risk ahead. 
%\noindent -Warning vs behavior support \\
%-Ghost vehicle as in game \\

\subsection{Short-Term Warning Elements}
\label{subsec:warning}
In order to emphasize the criticality of the situation, we propose to add further intuitive warning elements as e.g. pop-up signs and lane colorings. 
The following elements augment the proactive elements.

\subsubsection{Pop-up Signs}
Explicit symbols indicate the risk cause accompanied with the event time for collisions ($s_E$), distances to the risk spot for turns (i.e., right curve with $d_r$ and left curve with $d_l$) or stopping distance for crosswalks ($d_c$). In Fig. \ref{fig:popups}~a), the pop-up signs are pictured. 
% Besides the velocity difference, the risk type is an indication for the severity of the situation.
%Examples for collision risk are car-to-car crash., curve risk can be  as a single-car accident and regulatory risks will be a car-to-object collision. 
We want to stress that this is just a selection and more risk causes can be added. 
The purpose is also to clarify the reason for the warning and give more human-understandable information.

\subsubsection{Colored Events}
Finally, we highlight lane parts or positions according to the corresponding risks.  
% the determined color rating from the hazard route and velocity scale and relate the risks to the simulator environment. 
In the instance of curve and regulatory risk, the lane is colored from the ego position up to the point of maximal risk. 
For collision risk, we mark the point of the closest encounter as a red cube.
An illustration for regulatory risk induced from a stop line is depicted in Fig. \ref{fig:popups}~b). Again, the color is defined by the deviation $|v_0-v_{\text{tar}}|$. It also shows the therein considered navigation route with length $\Delta l_h$ and another unlikely path. 

It should be noted that the visualization of warnings only occurs if the risks are actually present. 
%\textcolor{red}{improve language, repeat intersection zone and navigation route}
%eingrauen unlikely paths and navigation path and describe in text, maybe delete Iz -> put line from unlikely path to green arrow
Altogether, the RNS provides a variety of tools to analyze and circumvent critical situations in intersection scenarios, while not overloading the driver's awareness.

\begin{figure}[t]
  \centering
  \resizebox{\linewidth}{!}{
  \import{img/}{colored_lane_new.pdf_tex}}  
  \vspace{-0.53cm}
  \caption{Short-term warning elements. Selected pop-up warnings (left) and colored lane (right).}
  \label{fig:popups} 
\end{figure} 


We present in section~\ref{ssec:faces} an application of PnP-HVAE on face images, using a pretrained state-of-the-art hierarchical VAE. 
Next, we study the application of our framework to natural images. To that end, we introduce  in section~\ref{ssec:patchVDVAE}  a patch hierachical VAE architecture, that is able to model natural images of different resolutions. In section~\ref{ssec:app_nat}, we provide deblurring, super-resolution and inpainting experiments to demonstrate the relevance of the proposed method.

Additional results are presented in Appendix~\ref{app:add}. All experiments can be reproduced using the code available at \url{https://github.com/jprost76/PnP-HVAE}.



\subsection{Face Image restoration (FFHQ)}\label{ssec:faces}
We first demonstrate the effectiveness of PnP-HVAE on highly structured data, by performing face image restoration.
Latent variable generative models can accurately model structured images such as face images \cite{karras2019style,vahdat2020nvae,child2021very,kingma2018glow}, and then be used to produce high quality restoration of such data. 
In our experiments, we use the VDVAE model of~\cite{child2021very}, pre-trained on the FFHQ dataset~\cite{karras2019style}, as our hierarchical VAE prior.
VDVAE has $L=66$ latent variable groups in its hierarchy and generates images at resolution $256\times256$.

We compare PnP-HVAE with the intermediate layer optimization algorithm (ILO)~\cite{daras2021intermediate} that is based on a different class of generative models than HVAE. ILO is a GAN inversion method which optimizes the image latent code along with the intermediate layer representation of a StyleGAN to generate an image consistent with a degraded observation.
We use the official implementation of ILO, along with a StyleGAN2 model~\cite{karras2020analyzing, stylegan2pytorch}, that was trained for 550k iterations on images of resolution $256\times256$ from FFHQ.  
As VDVAE and StyleGAN models are not trained on the same train-test split of FFHQ, we chose to evaluate the methods on a subset of 100 images from the CelebA dataset~\cite{liu2018large}. 
For super-resolution, the degradation model corresponds to the application of a gaussian low-pass filter followed by a $\times 4$ sub-sampling, and the addition of a gaussian white noise with $\sigma=3$.
For the deblurring, we considered motion blur and  gaussian kernels, both with a noise level $\sigma=8$. %

We provide quantitative comparisons in table~\ref{table:comp_ILO}, along with a visual comparison of the results in figure~\ref{fig:face_restoration}.
PnP-HVAE has the best  PSNR and SSIM results for all the considered restoration tasks, while ILO provides better results  for the perceptual distance.
By jointly optimizing the image and its latent variable, PnP-HVAE provides  results that are both realistic and consistent with the degraded observation.
On the other hand,  ILO  only optimizes on an extended latent space. This method generates  sharp and realistic images with better LPIPS scores,   
but the results lack  of consistency with respect to the observation, which explains the overall lower PSNR performance. 






\subsection{PatchVDVAE: a HVAE for natural images}\label{ssec:patchVDVAE}
Available generative models in the literature operate on images of  fixed resolutions and
are either restrained to datasets of limited diversity, or even to registered face images~\cite{kingma2018glow,child2021very, vahdat2020nvae, karras2019style}, or requiring additional class information~\cite{brock2018large, dhariwal2021diffusion, song2020score, luhman2022optimizing}.
Fitting an unconditional model on natural images appears to be a more difficult task, as their resolution can change, and their content is highly diverse.
The complexity of the problem can be reduced by learning a prior model on patches of reduced dimension. 
For image restoration problems, the patch model can be reused on images of higher dimensions~\cite{zoran2011learning,prost2021learning,altekruger2022patchnr}. When the model is a full CNN, the prior on the set of the  patches can  be computed efficiently by applying the network on the full image~\cite{prost2021learning}.

We thus introduce  patchVDVAE, a fully convolutional hierarchical VAE.
Contrary to existing HVAE models whose resolution is constrained by the constant tensor at the input of the top-down block, patchVDVAE can generate images of different resolutions by controlling the dimension of the input latent. 
This amounts to defining a prior on patches whose dimension corresponds to the receptive field of the VAE. A similar model is used for image denoising in~\cite{prakash2021interpretable}.

 
For PatchVDVAE architecture, we use the same bottom-up and top-down blocks as VDVAE~\cite{child2021very}, and replace the constant trainable input in the first top-down block by a latent variable, to make the model fully convolutional (details on the  architecture are given in Appendix~\ref{app:details}). 
The training dataset is composed of $128\times 128$ patches extracted from a combination of DIV2K~\cite{agustsson2017ntire} and Flickr2K~\cite{Lim_2017_CVPR_workshops} datasets.
We perform data augmentation by extracting  patches at $3$ resolutions: HR-images and $\times 2$ and $\times 4$ downscaled images. 
The model is trained for $7.10^5$ iterations with a batch size of $64$. Following the recommendation of~\cite{hazami2022efficient}, we use Adamax optimizer with an exponential moving average and gradient smoothing of the variance.
We set the decoder model to be a gaussian with diagonal covariance, as in~\cite{luhman2022optimizing}.
PatchVDVAE is fully convolutional and can generate images of dimension that are multiples of $64$ as illustrated by
figure~\ref{fig:vdvae}.

\newlength{\patchwidth}
\setlength{\patchwidth}{0.135\columnwidth}
\begin{figure}[!ht]
    \centering
    \begin{subfigure}[t]{.34\columnwidth}\hspace{0.1cm}
        \setlength{\tabcolsep}{0.02pt}
\renewcommand{\arraystretch}{0}
        \begin{tabular}{*{2}{p{1.03\patchwidth}}}
            \includegraphics[width=\patchwidth]{figures_arxiv/patchVDVAE/samples/generated/64x64/setup-5-image-0018.png} &
            \includegraphics[width=\patchwidth]{figures_arxiv/patchVDVAE/samples/generated/64x64/setup-5-image-0016.png} \\
            \includegraphics[width=\patchwidth]{figures_arxiv/patchVDVAE/samples/generated/64x64/setup-5-image-0008.png} &
            \includegraphics[width=\patchwidth]{figures_arxiv/patchVDVAE/samples/generated/64x64/setup-5-image-0019.png}   
        \end{tabular}
    \end{subfigure}\hspace{-0.15cm}
    \begin{subfigure}[t]{.64\columnwidth}
\begin{tabular}{cc}\vspace{-0.1cm}
\includegraphics[width=2\patchwidth]{figures_arxiv/patchVDVAE/samples/generated/256x256/setup-2-image-0009.png}&
        \includegraphics[width=2\patchwidth]{figures_arxiv/patchVDVAE/samples/generated/256x256/setup-2-image-0002.png}\end{tabular}

    \end{subfigure}
    \caption{\label{fig:vdvae} Left: $64\times64$ patches samples from our patchVDVAE model trained on patches from natural images.
    Right: PatchVDVAE is fully convolutional and it can generate images of higher resolution (here: $128\times128$).\vspace{-0.2cm}}
\end{figure}

\subsection{Natural images restoration}\label{ssec:app_nat}
We  evaluate PnP-HVAE on natural image restoration.
For each task, we report the average value of the PSNR, the SSIM, and the LPIPS metrics on $20$ images from the test set of the BSD dataset~\cite{MartinFTM01}.\\


\noindent
{\bf Image deblurring.}
In the experiments, we consider $2$ gaussian kernels and $2$ motion blur kernels from~\cite{levin2009understanding}, with $3$ different noise levels 
$\sigma \in \{2.55, 7.65, 12.75\}$.
As a baseline we consider  EPLL~\cite{zoran2011learning}, which learns a prior on image patches with a gaussian mixture model.
We also compare PnP-HVAE  with PnP-MMO and GS-PnP, $2$ competing convergent Plug-and-Play methods based on CNN denoisers.
PnP-MMO~\cite{pesquet2021learning} restricts the denoiser to be contraction in order to guarantee the convergence of the PnP forward-backard algorithm. GS-PnP~\cite{hurault2022gradient} considers a gradient step denoiser and reaches state-of-the-art performances of non converging methods~\cite{zhang2021plug}.
We set the temperature $\tau$  in our method as $0.95$, $0.8$ and $0.6$ for noise levels $2.55$, $7.65$ and $12.75$ respectively, and we let it run for a maximum of $50$ iterations. 
For the three compared methods we use the official implementations and pre-trained models provided by the respective authors. 
Details on the choice of hyperparameters for the concurrent methods are provided in the Appendix~\ref{app:details}
Figure~\ref{fig:deblurring_bsd} illustrates that our method provides correct deblurring results. 

According to table~\ref{tab:deb}, the performance of PnP-HVAE is between those of EPLL and GS-PnP and it outperforms PnP-MMO for large noise levels.\\

\begin{table}
\begin{center}\footnotesize
    \begin{tabular}{>{\centering}m{.3cm}*{5}{c}}
    $\sigma$ &Method & PSNR$\uparrow$ & SSIM$\uparrow$ & LPIPS$\downarrow$  \\ 
    \hline
    \multirow{4}{*}{\vcell{$2.55$}}
    & PnP-HVAE & $27.75$ & $0.79$ & $0.31$\\
    & GS-PNP \cite{hurault2022gradient} & $\mathbf{29.59}$ & $\mathbf{0.84}$ & $\mathbf{0.22}$\\
    & EPLL \cite{zoran2011learning} & $26.49$ & $0.71$ & $0.36$\\ 
    & PnP-MMO \cite{pesquet2021learning} & $\underbar{29.50}$ & $\underbar{0.83}$ & $\underbar{0.20}$ \\ \hline
    \multirow{4}{*}{\vcell{$7.65$}}
    & PnP-HVAE & $\underbar{26.36}$ & $\underbar{0.72}$ & $\underbar{0.40}$\\
    & GS-PNP \cite{hurault2022gradient} & $\mathbf{27.33}$ & $\mathbf{0.77}$ & $\mathbf{0.31}$\\
    & EPLL \cite{zoran2011learning} & $24.04$ & $0.66$ & $0.45$ \\ 
    & PnP-MMO \cite{pesquet2021learning} & $25.34$ & $0.69$ & $0.34$\\
    \hline
    \multirow{4}{*}{\vcell{$12.75$}}
    & PnP-HVAE & $\underbar{25.12}$ & $\mathbf{0.73}$ & $\underbar{0.47}$\\
    & GS-PNP \cite{hurault2022gradient} & $\mathbf{26.32}$ & $\mathbf{0.73}$ & $\mathbf{0.37}$\\
    & EPLL \cite{zoran2011learning} & $23.28$ & $0.61$ & $0.51$ \\ 
    & PnP-MMO \cite{pesquet2021learning} & $22.42$ & $0.53$& $0.54$ \\
    \hline
    &\vspace*{-.3cm}\\
            \multicolumn{2}{c}{Blur and motion kernels}& \multicolumn{3}{c}{
        \includegraphics*[scale=1]{figures_arxiv/kernels/4.png}\;\includegraphics*[scale=1]{figures_arxiv/kernels/7.png}\;\includegraphics*[scale=1]{figures_arxiv/kernels/9.png}\;\includegraphics*[scale=1]{figures_arxiv/kernels/11.png}} 
    \end{tabular}
        \caption{\label{tab:deb}Comparison  of PnP-HVAE  and other restoration methods on deblurring. Results are averaged on $4$ kernels.\vspace{-0.2cm}}% on image deblurring.}
    \end{center}
\end{table}

\begin{figure}
    
    \begin{subfigure}[h]{\linewidth}
        \centering
        \includegraphics*[width=\columnwidth]{figures_arxiv/deb_s255_k7.pdf}\vspace{-0.1cm}
        \caption{Gaussian blur, $\sigma=2.55$}
    \end{subfigure}
    \begin{subfigure}[h]{\linewidth}
        \centering
        \includegraphics*[width=\columnwidth]{figures_arxiv/deb_s765_k11.pdf}\vspace{-0.1cm}
        \caption{Motion blur, $\sigma=7.65$}
    \end{subfigure}\vspace*{-0.1cm}
    \caption{\label{fig:deblurring_bsd} Natural image deblurring\vspace{-0.1cm}}
\end{figure}

\noindent {\bf Effect of the temperature.}
PnP-HVAE gives control on the temperature of the prior over the latent space.
In figure~\ref{fig:temp_effect}, we illustrate that reducing the temperature increases the strength of the regularization prior. In this example the tuning $\tau=0.7$ produces the best performance.\\
\begin{figure}[!ht]
   
    \includegraphics[width=\columnwidth]{figures_arxiv/demo_temp.pdf}\vspace{-0.15cm}
    \caption{ \label{fig:temp_effect} Effect of the temperature in PnP-VAE on a deblurring problem, with $\sigma=7.65$.\vspace{-0.15cm}}
\end{figure}


\noindent
{\bf Image inpainting.}
Next we consider the task of noisy image inpainting. 
We compose a test-set of 10 images from the validation set of BSD~\cite{MartinFTM01} and we create masks
  by occluding diverse objects of small size in the images. 
A gaussian white noise with $\sigma=3$ is added to the images.
As a comparaison, we still consider GS-PnP and EPLL.
For PnP-HVAE, the temperature is set to $\tau=0.6$, and the algorithm is run for a maximum of $200$ iterations, unless the residual $||\x_{k+1}-\x_k||$ is on a plateau.
We provide on Table~\ref{tab:inpainting_bsd} the distortion metrics with the ground truth, as well as a visual
\begin{table}



\begin{center}
    \begin{tabular}{cccc}
        & PSNR$\uparrow$ & SSIM$\uparrow$ &LPIPS$\downarrow$ \\\hline
        PnP-HVAE  & $\mathbf{29.54}$ & $\mathbf{0.93}$ & $\mathbf{0.06}$\\
        GS-PNP & $28.52$ & $\mathbf{0.93}$ & $0.09$\\
        EPLL & $\underline{29.16}$ & $\mathbf{0.93}$ & $\mathbf{0.06}$\\
    \end{tabular}
    \caption{\label{tab:inpainting_bsd}Quantitative evaluation for inpainting on BSD.}
    \end{center}
\end{table}
comparison on figure~\ref{fig:inpainting_bsd}. 
With its hierarchical structure,  PnP-HVAE outperforms the compared methods. \vspace{0.05cm}



\begin{figure}[!h]
    \includegraphics[width=\columnwidth]{figures_arxiv/demo_inp_bsd2.pdf}\vspace{-0.1cm}
    \caption{\label{fig:inpainting_bsd}Natural image inpainting\vspace{-0.3cm}}
\end{figure}












%%%%%%%%%%%%%%%%
\section{Outlook}
\label{sec:Out}
%%%%%%%%%%%%%%%%

In this review we have attempted to present the main ideas behind the hypothesis
that the applicability of fluid dynamics to early phases of QGP dynamics can be
explained by a far-from-equilibrium hydrodynamic attractor. We have emphasised
the role of the kinematic setting specific to heavy-ion collisions, which makes
it plausible that such an attractor occurs also in QCD.  At the conceptual level
this picture is rather compelling. However, Bjorken flow is a very restrictive
setting, and despite some existing applications using the attractor in practical
calculations of phenomenologically interesting observables requires developing
effective methods of working with models with a large number of degrees of
freedom. Progress is likely to come gradually, by learning how to deal with
models of increasing complexity, extending the studies reviewed in Sections
\ref{sec:PhaseSpace} and \ref{sec:beyond}. 

The idea of hydrodynamic attractors has been closely connected with the
divergence of the hydrodynamic gradient expansion. In this review we have not
discussed this connection beyond its utilitarian aspects, but on a conceptual
level there have been some important developments in recent times. This includes
a proof of the generic divergence of the gradient expansion at the linearised
level without any symmetry assumptions, and its connection with the properties
of dispersion relations~\cite{Withers:2018srf,Heller:2020uuy}.  At the nonlinear
level, some of the results for Bjorken flow have been generalised to a much
wider class of flows, called longitudinal flows. In particular, the gradient
expansion has been shown to diverge for this class of
solutions~\cite{Heller:2021oxl}. It was also found that the large order
behaviour of the gradient series can be expressed in terms of new degrees of
freedom, the singulant fields, which track transient
effects~\cite{Heller:2021yjh}. The relevance of these advances to the study of
attractors remains an interesting challenge for the future. 

A number of issues were not addressed in this review. One is the presence of
other degrees of freedom, such as those  connected with chiral symmetry
breaking, and their possible effect on the attractor dynamics of
QGP~\cite{Mitra:2020mei,Mitra:2020hbj} (see also Ref.~\cite{Mitra:2022uhv}).
Another such issue is the role of fluctuations, which has been mostly neglected
in the attractor literature, with the notable exception of
Refs.~\cite{Akamatsu:2016llw,Chen:2022ryi}. 

It would also be very interesting to understand the role of quantum effects in
the early-time dynamics. Of course the hydrodynamic picture implicitly contains
them, but in a rather opaque way. On the other hand, the kinetic theory
description arises from quantum field theory through a series of
approximations~\cite{Mueller:2002gd,Jeon:2004dh} and it should be possible to
study the origin and robustness of the kinetic theory attractor in a framework
which allows for a systematic incorporation of quantum corrections. This is
connected with other approaches to early-time dynamics, including those
involving ideas such as the Color Glass Condensate or non-thermal
attractors~\cite{Berges:2008wm,Mazeliauskas:2018yef,Brewer:2019oha,Brewer:2022ifw,Brewer:2022vkq,Berges:2020fwq}.
It is not yet fully understood how they are related to the ideas reviewed here,
and clarifying this appears to be a very promising avenue for future research. 

Finally, there is the more general question about far-from-equilibrium
attractors in nonequilibrium systems.  In the context of heavy-ion collisions
the specific kinematic circumstances have lead us to consider boost-invariant
expansion where far-from-equilibrium hydrodynamic behaviour was first noted, but
there could be other situations where analogous phenomena might
appear~\cite{Baggioli:2021tzr}, perhaps even in the non-relativistic
domain~\cite{Le:2022ntg}.  Another context where far from equilibrium
hydrodynamic attractors can occur is the dynamics of systems in nontrivial
spacetime backgrounds (see e.g. Ref.~\cite{Vyas:2022hkm}). 





% use section* for acknowledgment
\section*{Acknowledgment}
\noindent This work has been supported by the European Unions Horizon 2020 project \textit{VI-DAS}, under the grant agreement number 690772. 
Map data \copyright~OpenStreetMap contributors, licensed under the Open Database License. 
We like to thank Cristhiam Felipe Pulido Riveros from National University of Colombia for his support. %on mainting the local dynamic map.
\vspace{0.15cm}

\bibliographystyle{IEEEtran}
%{\footnotesize \bibliography{bib.bib}}
\bibliography{bib}

% that's all folks
\end{document}


